\documentclass{article}
\title{Recitation notes}
\author{Jordan Hoffart}
\date{August 20, 2024}
\usepackage{amsmath,amsthm,amssymb}
\theoremstyle{plain}
\newtheorem{theorem}{Theorem}
\begin{document}
\theoremstyle{definition}
\newtheorem{definition}{Definition}
\maketitle
\section{Fundamental theorem of calculus}
\begin{definition}[Antiderivative]
	Let $f$ be a function defined on an interval $I$.
	An antiderivative of $f$ is a differentiable function $F$ defined on $I$ such that
	\begin{equation}
		F'(x) = f(x)
	\end{equation}
	for all $x$ in $I$.
\end{definition}
\begin{theorem}[Fundamental theorem of calculus, part 1]
	If $f$ is a continuous function on an interval $[a,b]$ and we define the function $F$ on $[a,b]$ by
	\begin{equation}
		F(x) = \int_a^xf(t)\,dt
	\end{equation}
	then $F$ is an antiderivative of $f$ on $[a,b]$.
\end{theorem}
\begin{theorem}[Fundamental theorem of calculus, part 2]
	If $f$ is a continuous function on an interval $[a,b]$
	and if $F$ is an antiderivative of $f$ on $[a,b]$, then
	\begin{equation}
		\int_a^b f(x)\,dx = F(b) - F(a).
	\end{equation}
\end{theorem}
\section{Substitution rule}
\begin{definition}[Indefinite integral]
	If $f$ is a function with an antiderivative $F$, then the indefinite integral of $f$ is the collection of all antiderivatives of $f$.
	We denote this collection by 
	\begin{equation}
		\int f(x)\,dx.
	\end{equation}
	The use of the letter $x$ in our notation is arbitrary.
	We can use any other symbol, as long as we are consistent.
	That is, all of these notations represent the indefinite integral of $f$:
	\begin{equation}
		\int f(x)\,dx = \int f(y)\,dy = \int f(z)\,dz = \int f(u)\,du = \dots
	\end{equation}
	as well as any other choices of the symbol of integration.
\end{definition}
\begin{theorem}[Substitution rule for indefinite integrals]
	If $g$ is a differentiable and invertible function, and if $f$ is a continuous function on the range of $g$, then 
	\begin{enumerate}
		\item For every antiderivative $F$ of $f$, $F \circ g$ is an antiderivative of $(f \circ g)g'$.
		\item For every antiderivative $G$ of $(f\circ g)g'$, $G \circ g^{-1}$ is an antiderivative of $f$.
	\end{enumerate}
	We summarize this by writing
	\begin{equation}
		\int f(g(x))g'(x)\,dx = \int f(u)\,du
	\end{equation}
	with $u = g(x)$.
\end{theorem}
\begin{proof}
	Let $F$ be an antiderivate of $f$.
	Then $(F\circ g)'(x) = F'(g(x))g'(x) = f(g(x))g'(x)$, so $F \circ g$ is an antiderivative of $(f \circ g)g'$.
	
	Now let $G$ be an antiderivative of $(f\circ g)g'$.
	Then 
	\[(G \circ g^{-1})'(x) = G'(g^{-1}(x))\frac{1}{g'(g^{-1}(x))} = f(x),\]
	so $G \circ g^{-1}$ is an antiderivative of $f$.
\end{proof}
\begin{theorem}[Substitution rule for definite integrals]
	If $g$ is a differentiable and invertible function on an interval $[a,b]$, and if $f$ is a continuous function on the interval between $g(a)$ and $g(b)$, then 
	\begin{equation}
		\int_a^b f(g(x))g'(x)\,dx = \int_{g(a)}^{g(b)}f(u)\,du.
	\end{equation}
\end{theorem}
\end{document}
